\documentclass[12pt,a4paper]{article}

% Packages
\usepackage[margin=1in]{geometry}
\usepackage{amsmath, amssymb, amsthm}
\usepackage{graphicx}
\usepackage{bm}
\usepackage{hyperref}
\usepackage{titlesec}

% Formatting
\titleformat{\section}{\large\bfseries}{\thesection}{1em}{}
\titleformat{\subsection}{\normalsize\bfseries}{\thesubsection}{1em}{}

% Title
\title{\textbf{Column Buckling and Effective Length in Steel Structures}}
\author{Prateeksha Sharma}
\date{}

\begin{document}

\maketitle

\begin{abstract}
Column buckling is a fundamental problem in steel structures, where a compression member can fail at loads lower than its yield capacity due to instability. This article discusses Euler's buckling theory and the concept of effective length.
\end{abstract}

\section{Introduction}
Columns are compressive load carrying members in steel structures. The faiklure modes in columns can be material yielding for stocky columns while the slender columns often fail by bucknling, a sudden lateral deflection under compressive axial load.

\section{Euler's Buckling Theory}
For an ideal, perfectly straight column with pinned ends and subjected to a concentric axial load, Leonhard Euler derived the critical buckling load as:
\begin{equation}
    P_{cr} = \frac{\pi^2 E I}{(K L)^2}
\end{equation}
where:
\begin{itemize}
    \item $E$ = modulus of elasticity,
    \item $I$ = least moment of inertia,
    \item $L$ = actual column length,
    \item $K$ = effective length factor.
\end{itemize}

\subsection{Slenderness Ratio}
The slenderness ratio is defined as:
\begin{equation}
    \lambda = \frac{K L}{r}
\end{equation}
where $r = \sqrt{I/A}$ is the radius of gyration. Columns with high $\lambda$ values are more prone to buckling.

\section{Effective Length Factor}
The effective length factor $K$ accounts for end boundary conditions:
\begin{itemize}
    \item Pinned–Pinned: $K = 1.0$
    \item Fixed–Fixed: $K \approx 0.5$
    \item Fixed–Free (cantilever): $K = 2.0$
    \item Fixed–Pinned: $K \approx 0.7$
\end{itemize}
In real structures, $K$ is influenced by frame stiffness, member continuity, and joint rigidity.

\section{Inelastic Buckling}
When $\lambda$ is moderate, yielding may occur before elastic buckling. The inelastic buckling capacity can be estimated using empirical design curves.

\section{LRFD Design Procedure}
According to AISC 360-16, the nominal compressive strength $P_n$ is:
\begin{equation}
    P_n =
    \begin{cases}
        F_{cr} A_g & \text{if elastic buckling governs}, \\
        F_{cr} A_g & \text{if inelastic buckling governs},
    \end{cases}
\end{equation}
where:
\[
F_{cr} = 
\begin{cases}
0.658^{F_y / F_e} F_y & \text{if } \lambda_c \leq 1.5, \\
0.877 F_e & \text{if } \lambda_c > 1.5
\end{cases}
\]
and $F_e = \frac{\pi^2 E}{(K L / r)^2}$ is the elastic critical stress.

\section{Conclusion}
Effective length factor K is essential for accurate column design as it affects the buckling capacity. It is influenced by frame stability, joint stiffness and end conditions. Accurate K values result in efficient designs whereas conservative assumptions ensure more safety.
\section*{References}
\begin{enumerate}
    \item AISC 360-16, \textit{Specification for Structural Steel Buildings}.
    \item Geschwindner, L. F. (2012). Unified design of steel structures (2nd ed.). John Wiley & Sons.
    
\end{enumerate}

\end{document}
